\chapter*{Introducción}
\addcontentsline{toc}{chapter}{Introducción}

Este libro reúne ejercicios de LeetCode frecuentes en entrevistas de Amazon,
resueltos en C++ y explicados para todo público. Cada encabezado de problema
es un enlace: al hacer clic se abre el archivo \texttt{.cpp} correspondiente
en Visual Studio Code, donde el código es compilable y trae casos de prueba.

\paragraph{Dónde está el código.} Todo el código fuente vive en la carpeta
\texttt{src/} del repositorio, organizada por bloque (\texttt{src/01-grafos/},
\texttt{src/02-arboles/}, \ldots). El libro lo importa directamente desde ahí
(no se duplica). Bajo el título de cada problema se muestra la ruta exacta del
archivo, por si quieres abrirlo o consultarlo directamente en \texttt{src/}.

\paragraph{Animaciones interactivas.} Además del código, cada problema incluye
el enlace \emph{\ensuremath{\triangleright}\ Ver animación} con dos accesos:
\textbf{local} (abre el visualizador del repo descargado, funciona sin internet)
y \textbf{web} (abre la versión publicada en GitHub Pages desde cualquier equipo).
Muestra una visualización del algoritmo paso a paso: la cuadrícula se pinta, la
cola del BFS avanza, el árbol se recorre. El visualizador vive en la carpeta
\texttt{rep-visual/}, funciona sin instalar nada (basta un navegador) y permite
alternar el idioma entre español e inglés. Los enlaces piden abrirse en una
ventana nueva para no reemplazar la vista del PDF.

\paragraph{Cómo compilar.} Ejecuta \texttt{build.ps1} en PowerShell. El script
detecta la ruta del repositorio e inyecta las rutas de los enlaces, luego
compila con \texttt{latexmk -xelatex}.
